\specialchapter{引言}

在本书第 III 章至第 V 章中，我们详细考察实或复拓扑向量空间的自同态代数。在有限维空间的情形，这一代数的元素逐个来看的主要性质都可由若尔当分解（A, VII, p. 42, § 5, 第 9 目）导出。我们在此研究它们向无穷维空间的推广，然而这推广要更为精细。

在第 III 章中，我们在非常一般的框架下考察拓扑向量空间之间的紧线性映射。它们的谱性质与有限维情形最为接近。例如，在巴拿赫空间的特殊情形，紧自同态的谱中的非零元素都是具有有限谱重数的广义特征值。

在希尔伯特空间的框架下，会出现一种易于处理的分类和结构。这正是第 IV 章的对象。该章从希尔伯特空间中紧算子的研究开始：它们借助奇异值的描述，以及正规自同态情形下的对角化，都完全类似于有限维空间上已知的性质。更一般地，谱定理给出了希尔伯特空间正规自同态的简单模型。

某些重要应用，例如涉及微分算子的应用，但也涉及量子力学的应用，要求一种更精微的理论：即定义在希尔伯特空间的稠密子空间上的线性映射---未必连续---的理论，我们把这类映射称为部分算子。在这一框架下，自伴算子或正规算子的定义本身就必须极其精确，否则便会产生重大的谬误。一旦掌握了这些定义，就会出现一种与连续自同态情形相类似的分类；而对有界或无界的普遍可测函数的函数演算，则在连续正规自同态与正规部分算子之间建立起紧密的联系。

顺带地，我们还会谈到带迹的正规自同态，特别是希尔伯特空间 \(L^{2}(X,\mu)\) 中带积分核的积分自同态。人们熟知，在自守形式与表示的理论中，断言这类自同态（当其存在时）的迹等于核沿对角线的积分的那条公式具有至关重要的意义。我们还用一节的篇幅叙述 \(R^{n}\) 中广义函数与缓增广义函数的基本定义和性质。这些概念不仅给出了欧几里得傅里叶理论诸基本原理（参见 II, § 1）的最方便的表述，而且还使人们能够确定许多重要微分算子的定义域。

最后，在第 V 章中，我们把希尔伯特空间自同态的研究与群的酉表示联系起来。在若干初等的一般性考虑之后---其中舒尔引理是最重要的命题---我们研究局部紧群的情形。然后我们考察正定函数，并给出它们的最初应用。这些应用之一是盖尔范德--奈马克--塞加尔定理，它断言每个星代数都同构于某个希尔伯特空间的自同态代数的一个闭子代数。另一些应用则涉及局部紧群的不可约酉表示所成的空间的基本性质，特别是盖尔范德--赖科夫定理，它确立了这些表示分离点。在本章最后一节中，我们证明关于紧群表示的基本结果，它们曾在《群与李代数》第 IX 章中为我们所用。

A. BOREL, Representations of locally compact groups, Lecture Notes in Mathematics 276, Springer, 1972.

M. REED 与 B. SIMON, Methods of Mathematical Physics, Academic Press, Vol. 1, 1980 与 Vol. 2, 1975.

K. SCHMÜDGEN, Unbounded Self-adjoint Operators on Hilbert Space, Graduate Texts in Mathematics 265, Springer, 2012.

在前面的各章与各书中，对“TS, to appear”的下列引用应作如下修改：

EVT, V, p. 35, 第 1 行。应将“TS, IV, § 6”改为“(TS, IV, p. 160, remark)”。

INT, IX, p. 94, § 6, 第 12 目。应将“One can prove a converse \ldots{} bounded”改为“The converse of prop. 11 is a particular case of TS, V, p. 455, th. 5.”

LIE, IX, p. 61, § 6, 第 4 目, 第 6 行。应将“TS, to appear”改为“TS, III, p. 36, prop. 9”。

LIE, IX, p. 63, § 6, 第 4 目, 第 6 行。应将“TS, to appear”改为“TS, III, p. 36, rem.”。

LIE, IX, p. 66, § 7, 注 1。应将“to a chapter to appear”改为“to paragraph 4 of Chapter V”。

LIE, IX, p. 69, § 7, 第 2 目。应将“then according to TS”改为“then (TS, V, p. 464, prop. 6)”。

LIE, IX, p. 71, § 7, 第 3 目。应将“which, according to TS, is an isomorphism”改为“which, according to TS, V, p. 469, cor. 2, is an isomorphism.”

LIE, IX, p. 72, § 7, 第 4 目。应将“(TS)”改为“(TS, V, p. 459, cor. 3)”。

LIE, IX, p. 74, § 7, 第 4 目。在“of the space of central representative functions on G”之后加上“(TS, V, p. 469, cor. 2, a)”；在“of the space ZL²(G) of central square-integrable functions on G”之后加上“(TS, V, p. 468, prop. 9).”

LIE, IX, p. 78, § 8, 第 1 目。应将“In this number, one recalls definitions and results from TS”改为“One recalls definitions and results from TS, V, p. 386, no. 8; p. 392, no. 10; p. 470, no. 7.”

LIE, IX, p. 84, § 8, 第 3 目。应将“Recall (TS) the formulas”改为“One has (TS, V, p. 467, formulas (1) and (2)) the formulas”。

LIE, IX, p. 89, § 9, 第 1 目。应将“TS, III, §2, no. 7, prop. 16”改为“TS, III, p. 67, prop. 10”。

LIE, IX, p. 90, § 9, 第 1 目。应将“TS, III, §2, no. 6”改为“TS, III, p. 61, no. 4”。

LIE, IX, p. 91, § 9, 第 2 目。应将“TS, to appear”改为“TS, III, p. 36, prop. 9 and TS, V, p. 464, prop. 7”。

LIE, IX, p. 92, 注 1。应将“(TS, to appear)”改为“(TS, V, p. 379, def. 3)”。

LIE, IX, p. 99, 附录 1, 第 1 目。应将“(TS, to appear)”改为“(TS, V, p. 464, remark 4)”。

LIE, IX, p. 128, § 9 习题 3 (b)。应将“(TS, to appear)”改为“(Use the Peter--Weyl theorem, TS, V, p. 462, th. 1)”。

TS, I, p. 103, 例 6。应将“V, to appear”改为“V, p. 442, th. 2.”

TS, I, p. 173, 习题 5, d)。应将“III, to appear”改为“III, p. 82, § 6.”

TS, II, p. 237, 第 9 目。应将“IV, to appear”改为“IV, p. 196, § 3”。

TS, II, p. 251, 第 1 目。应将“V, to appear”改为“V, p. 425, prop. 10”。

\chapter{紧线性映射与扰动}
\setcounter{exercisepar}{0}

在本章中，所考虑的向量空间都是等于 R 或 C 的域 K 上的向量空间。本书中所作的关于 EVT, II 的引用一般针对 K = R 的情形；关于 K = C 的情形，参见 EVT, II, p. 64 至 68。

半赋范空间是指赋有一个半范数 p 以及由 p 定义的拓扑的向量空间 E（参见 EVT, II, p. 2）。E 的单位球，或 p 的单位球，是 E 中满足 \(p(x) \leqslant 1\) 的元素 x 所成的集。

对每个向量空间 E，我们用 \(1_{E}\) 表示 E 的恒等映射。若 E, F 和 G 是向量空间，且 u: \(E \rightarrow F\)、\(v: F \rightarrow G\) 是线性映射，我们有时用 vu 表示从 E 到 G 的线性映射 \(v \circ u\)。

设 E 和 F 是拓扑向量空间。从 E 到 F 的连续线性映射也称为态射，而当 F = E 时称为自同态。从 E 到 F 的双射线性映射，若它连续且其逆也连续，则称为同构，而当 F = E 时称为自同构。当 E 和 F 是半赋范空间时，自同态、同构和自同构这几个概念都是就底层的拓扑向量空间结构而言的。

给定拓扑向量空间 E 和 F，我们用 \(\mathcal{L}^{\mathrm{f}}(\mathrm{E};\mathrm{F})\) 表示从 E 到 F 的有限秩连续线性映射所成的向量空间。我们也用 \(\mathcal{L}^{\mathrm{f}}(\mathrm{E})\) 表示向量空间 \(\mathcal{L}^{\mathrm{f}}(\mathrm{E};\mathrm{E})\)。

我们回顾，从 E 到 F 的连续线性映射 u 称为严格的，如果它通过过渡到商空间诱导出从 E/ Ker(u) 到 u(E) 上的一个同构（TG, III, p. 16, def. 1）；这等价于说，E 中 0 的每个邻域的像是 u(E) 中 0 的一个邻域（TG, III, p. 16, prop. 24）。

\section{紧线性映射}

\subsection{紧线性映射}

定义 1. --- 设 E 是拓扑向量空间，F 是分离的拓扑向量空间。从 E 到 F 的线性映射 u 称为紧的，如果存在 E 中 0 的一个邻域 V，使得 \(u(V)\) 是 F 的一个相对紧子集。

我们用 \(\mathcal{L}^{\mathrm{c}}(\mathrm{E};\mathrm{F})\) 表示从 E 到 F 的紧线性映射所成的集；我们也用 \(\mathcal{L}^{\mathrm{c}}(\mathrm{F})\) 表示集 \(\mathcal{L}^{\mathrm{c}}(\mathrm{F};\mathrm{F})\)。

注记. --- 1) E 中 0 的一个邻域吸收 E 的每个有界子集；因此，E 的有界子集在从 E 到 F 的紧线性映射下的像是 F 的一个相对紧子集。

\begin{enumerate}
\def\labelenumi{\arabic{enumi})}
\setcounter{enumi}{1}
\item
  设 E 是半赋范空间，B 是 E 的单位球。为使从 E 到一个分离拓扑向量空间 F 的线性映射 u 是紧的，必须且只需 \(u(\mathrm{B})\) 是 F 的一个相对紧子集。
\item
  设 F 是分离的拓扑向量空间。为使 F 的向量子空间 E 是有限维的，必须且只需 E 到 F 中的典范单射是紧的（EVT, I, p. 15, th. 3）。
\item
  对域 R 上的每个拓扑向量空间 E，我们用 \(\mathrm{E}_{(\mathbf{C})}\) 表示 E 的复化拓扑向量空间（EVT, II, p. 65）。我们把 E 等同于 \(\mathrm{E}_{(\mathbf{C})}\) 的一个实拓扑向量子空间（经由映射 \(x \mapsto 1 \otimes x\)）；于是 \(\mathrm{E}_{(\mathbf{C})}\) 的底层实拓扑向量空间是 E 与 iE 的拓扑直和。
\end{enumerate}

设 \(u\) 是从 E 到 \(\mathbf{R}\) 上的一个分离拓扑向量空间 F 的线性映射。用 \(u_{(\mathbf{C})}\) 表示 \(\mathbf{C}\)-线性映射（它从 \(\mathrm{E}_{(\mathbf{C})}\) 到 \(\mathrm{F}_{(\mathbf{C})}\) 中，并且拓广 \(u\)）。形如 \(V + iV\) 的集合（其中 \(V\) 是 E 中 0 的一个邻域）构成 \(\mathrm{E}_{(\mathbf{C})}\) 中 0 的一个基本邻域系。为使集合 \(u_{(\mathbf{C})}(V + iV) = u(V) + iu(V)\) 在 \(\mathrm{F}_{(\mathbf{C})}\) 中相对紧，必须且只需 \(u(V)\) 在 F 中相对紧。因此，为使 \(\mathbf{C}\)-线性映射 \(u_{(\mathbf{C})}\) 是紧的，必须且只需 \(\mathbf{R}\)-线性映射 \(u\) 是紧的。

出于同样的思路，当 \(v\) 是从 \(\mathrm{E}_{(\mathbf{C})}\) 到一个 C-拓扑向量空间 G 的连续 C-线性映射时，为使 \(v\) 是紧的，必须且只需它在 E 上的限制是紧的。

\begin{enumerate}
\def\labelenumi{\arabic{enumi})}
\setcounter{enumi}{4}
\item
  设 E 是拓扑向量空间，F 是分离的拓扑向量空间。设 \(F_{1}\) 是 F 的一个闭向量子空间，j 是 \(F_{1}\) 到 F 中的典范单射。为使从 E 到 \(F_{1}\) 的线性映射 u 是紧的，必须且只需从 E 到 F 的线性映射 \(j \circ u\) 是紧的。
\item
  设 \((\mathrm{G}_{i})_{i\in\mathrm{I}}\) 是一族分离的拓扑向量空间，且对每个 \(i\in I\)，设 \(u_{i}\) 是从拓扑向量空间 E 到 \(G_{i}\) 的一个紧线性映射。若集合 I 有限，或者 E 是半赋范空间，则线性映射 \(x\mapsto(u_{i}(x))_{i\in\mathrm{I}}\)（从 E 到 \(\prod_{i\in I}G_{i}\) 中）是紧的（注记 2 及 TG, I, p. 63, Th. 3）。
\end{enumerate}

命题 1. --- 设 E 是拓扑向量空间，F 是分离的拓扑向量空间。

\begin{enumerate}
\def\labelenumi{\alph{enumi})}
\item
  从 E 到 F 的每个紧线性映射都是连续的。
\item
  从 E 到 F 的每个有限秩连续线性映射都是紧的；
\item
  从 E 到 F 的紧线性映射所成的集 \(\mathcal{L}^{\mathrm{c}}(\mathrm{E};\mathrm{F})\) 是 \(\mathcal{L}(\mathrm{E};\mathrm{F})\) 的一个向量子空间。
\end{enumerate}

我们来证明 a)。设 \(v: E \rightarrow F\) 是紧线性映射，U 是 F 中 0 的一个邻域。取 E 中 0 的一个邻域 V，使得 \(v(V)\) 在 F 中相对紧。于是 \(v(V)\) 有界，因此存在实数 \(\lambda > 0\)，使得集合 \(v(\lambda V)\) 包含于 U；从而 v 连续。

设 \(u \colon \mathrm{E} \to \mathrm{F}\) 是有限秩连续线性映射。由于 \(u(\mathrm{E})\)（即 \(u\) 的像）是分离的且是有限维的，故存在 \(u(\mathrm{E})\) 中 0 的一个紧邻域 A。集合 \(\mathrm{V} = u^{-1}(\mathrm{A})\) 是 E 中 0 的一个邻域，并且 \(u(\mathrm{V}) \subset \mathrm{A}\)，所以线性映射 \(u\) 是紧的。这就证明了断言 \(b\)）。

设 \(u_{1}\) 和 \(u_{2}\) 是从 E 到 F 的紧线性映射。取 E 中 0 的邻域 \(V_{1}\) 和 \(V_{2}\)，使得集合 \(u_{1}(V_{1})\) 与 \(u_{2}(V_{2})\) 都在 F 中相对紧。令 \(V = V_{1} \cap V_{2}\)，\(u = u_{1} + u_{2}\)。于是 \(u(V)\) 包含于 \(u_{1}(V_{1}) + u_{2}(V_{2})\)，从而在 F 中相对紧，故线性映射 u 是紧的。\(\mathcal{L}^{c}(E;F)\) 在位似变换下稳定是显然的。这就证明了 \(\mathcal{L}^{c}(E;F)\) 是 \(\mathcal{L}(E;F)\) 的一个向量子空间，从而得到断言 c)。

命题 2. --- 设 E 是半赋范空间，F 是分离的拓扑向量空间。假设 F 的每个有界闭子集都是完备的。那么向量子空间 \(\mathcal{L}^{\mathrm{c}}(\mathrm{E};\mathrm{F})\)（\(\mathcal{L}(\mathrm{E};\mathrm{F})\) 的）对有界收敛拓扑是闭的。

设 \(u\) 是 \(\mathcal{L}(\mathrm{E};\mathrm{F})\) 中对有界收敛拓扑附着于 \(\mathcal{L}^{\mathrm{c}}(\mathrm{E};\mathrm{F})\) 的一个元素。用 B 表示 E 的单位球，V 表示 F 中 0 的一个邻域。取 F 中 0 的一个邻域 W，使得 \(\mathrm{W} + \mathrm{W} \subset \mathrm{V}\)。按假设，存在元素 \(v\)（它属于 \(\mathcal{L}^{\mathrm{c}}(\mathrm{E};\mathrm{F})\)），使得 \((u - v)(\mathrm{B}) \subset \mathrm{W}\)。集合 \(v(\mathrm{B})\) 在 F 中相对紧，故存在 F 的一个有限子集 M，使得 \(v(\mathrm{B}) \subset \mathrm{M} + \mathrm{W}\)，从而 \(u(\mathrm{B}) \subset \mathrm{M} + \mathrm{V}\)。由于这对每个 V 都成立，集合 \(u(\mathrm{B})\) 在 F 中预紧，因而其闭包也预紧（TG, II, prop. 1, p. 30）。后者是闭的且有界，故由对 F 所作的假设，它是完备的，从而是紧的（TG, II, cor., p. 30）。由此推出 \(u(\mathrm{B})\) 相对紧。这就证明了 \(u\) 属于 \(\mathcal{L}^{\mathrm{c}}(\mathrm{E};\mathrm{F})\)。

若 F 是局部凸的，则命题的假设意味着 F 是拟完备的（EVT, III, p. 8, def. 6）。

推论. --- 设 E 和 F 是巴拿赫空间。集合 \(\mathcal{L}^{\mathrm{c}}(\mathrm{E};\mathrm{F})\) 是巴拿赫空间 \(\mathcal{L}(\mathrm{E};\mathrm{F})\) 的一个闭向量子空间。它包含从 E 到 F 的有限秩连续线性映射所成的空间 \(\mathcal{L}^{\mathrm{f}}(\mathrm{E};\mathrm{F})\) 的闭包。

这由命题 1 和命题 2 得出。

\pandocbounded{\includegraphics[keepaspectratio]{images/part001_746f5095.jpg}}

存在巴拿赫空间 E 和 F，使得 \(\mathcal{L}^{\mathrm{f}}(\mathrm{E};\mathrm{F})\) 在 \(\mathcal{L}^{\mathrm{c}}(\mathrm{E};\mathrm{F})\) 中不稠密（参见 III, p. 16, 注记 6 及 III, p. 19, 定理 4, b)）。

命题 3. --- 设 \(\mathrm{E}_{1}\)、E、F、\(\mathrm{F}_{1}\) 是拓扑向量空间，其中 F 和 \(\mathrm{F}_{1}\) 假设为分离的。设 v: \(\mathrm{E}_{1}\) → E、u: E → F、w: F → \(\mathrm{F}_{1}\) 是线性映射。若 v 和 w 连续而 u 紧，则 w ∘ u ∘ v 紧。

按假设，存在 E 中 0 的一个邻域 V，使得 \(u(\mathrm{V})\) 在 F 中相对紧。令 \(\mathrm{U} = v^{-1}(\mathrm{V})\)。于是 U 是 \(\mathrm{E}_{1}\) 中 0 的一个邻域，且它在 \(w \circ u \circ v\) 下的像包含于 \(w(u(\mathrm{V}))\)，从而在 \(\mathrm{F}_{1}\) 中相对紧。因此，线性映射 \(w \circ u \circ v\) 是紧的。

设 E 是分离的拓扑向量空间；由命题 1 和命题 3，\(\mathcal{L}^{\mathrm{c}}(\mathrm{E})\) 是代数 \(\mathcal{L}(\mathrm{E})\) 的一个双边理想。当 E 是巴拿赫空间时，\(\mathcal{L}(\mathrm{E})\) 是巴拿赫代数，而 \(\mathcal{L}^{\mathrm{c}}(\mathrm{E})\) 是 \(\mathcal{L}(\mathrm{E})\) 的一个闭双边理想（命题 2 的推论）。若 E 是无穷维的，则它是一个真理想（参见 p. 2, 注记 3）。

推论. — 设 E 是希尔伯特空间。空间 \(\mathcal{L}^{\mathrm{c}}(\mathrm{E})\) 是 \(\mathcal{L}(\mathrm{E})\) 的一个闭的自伴双边理想。特别地，它是一个星代数。

事实上，空间 \(\mathcal{L}^{c}(E)\) 是 \(\mathcal{L}(E)\) 的一个闭双边理想，从而是自伴的（I, p. 122, 引理 15）。

命题 4. — 设 E 是拓扑向量空间，\(\widehat{\mathrm{E}}\) 是它的分离完备化，\(j\) 是 E 到 \(\widehat{\mathrm{E}}\) 中的典范映射。设 \(u\) 是从 E 到一个分离拓扑向量空间 F 的紧线性映射。则存在唯一的紧线性映射 \(v\)（从 \(\widehat{\mathrm{E}}\) 到 F 中），使得 \(u = v \circ j\)。

把 F 等同于 \(\widehat{\mathrm{F}}\) 的一个拓扑向量子空间，并用 \(\widehat{u}\colon \widehat{\mathrm{E}}\to \widehat{\mathrm{F}}\) 表示使 \(\widehat{u}\circ j\) 在 E 上与 \(u\) 重合的唯一连续线性映射。由于 \(u\) 紧，存在 E 中 0 的一个邻域 V 和 F 的一个紧子集 A，使得 \(u(\mathrm{V})\subset \mathrm{A}\)。我们有 \(\widehat{u}(j(\mathrm{V}))\subset \mathrm{A}\)，从而 \(\widehat{u} (\overline{j(\mathrm{V})})\subset \mathrm{A}\)。而 \(\overline{j(\mathrm{V})}\) 是 \(\widehat{\mathrm{E}}\) 中 0 的一个邻域（TG, III, p. 24, prop. 7）。因此，\(\widehat{u}\) 的像包含于 F，且连续线性映射 \(v\colon \widehat{\mathrm{E}}\rightarrow \mathrm{F}\)（由 \(\widehat{u}\) 过渡到子空间得到的）是紧的。由于 \(j(\mathrm{E})\) 在 \(\widehat{\mathrm{E}}\) 中稠密，映射 \(v\) 是从 \(\widehat{\mathrm{E}}\) 到 F 的满足 \(u = v\circ j\) 的唯一连续线性映射。

命题 5. --- 设 E 是局部凸空间，F 是分离的拓扑向量空间，u 是从 E 到 F 的紧线性映射。则存在一个巴拿赫空间 G、一个从 E 到 G 的连续线性映射 v 和一个从 G 到 F 的紧线性映射 w，使得 \(u = w \circ v\)。

取 E 中 0 的一个邻域 V，使得 \(u(\mathrm{V})\) 在 F 中相对紧。设 \(p\) 是 E 上的一个连续半范数，使得 V 包含 \(p\) 的单位球（EVT, II, p. 26, cor.）。用 \(\mathrm{E}_p\) 表示给 E 赋以半范数 \(p\) 所得到的半赋范空间。映射 \(u\) 是从 \(\mathrm{E}_p\) 到 F 的紧线性映射。\(\mathrm{E}_p\) 的分离完备化 G 是一个巴拿赫空间。用 \(v\) 表示从 \(\mathrm{E}_p\) 到 G 的典范线性映射。由命题 4，存在一个紧线性映射 \(w: \mathrm{G} \to \mathrm{F}\)，使得 \(u = w \circ v\)。另一方面，v 是从 E 到 G 的连续映射，因为 E 的拓扑比 \(\mathrm{E}_p\) 的拓扑更细。

\subsection{紧线性映射与弱拓扑}

设 E 是局部凸空间，\(E'\) 是它的对偶。回顾，E 上的拓扑 \(\sigma(\mathrm{E},\mathrm{E}')\) 称为 E 上的弱拓扑（TVS, IV, p. 4）。在本节中，\(E_{\sigma}\) 表示赋有弱拓扑的空间 E。

命题 6. --- 设 E 是局部凸空间，F 是分离的局部凸空间，\(u: \mathrm{E} \to \mathrm{F}\) 是紧线性映射，B 是 E 的一个有界子集。\(u\) 在 B 上的限制是从赋有 \(\sigma(\mathrm{E}, \mathrm{E}')\) 诱导拓扑的集合 B 到空间 F 的连续映射。

映射 \(u\) 是从 E 到 F 的连续线性映射，也是从 \(\mathrm{E}_{\sigma}\) 到 \(\mathrm{F}_{\sigma}\) 的连续线性映射。因此，它在 B 上的限制是从赋有 \(\sigma(\mathrm{E}, \mathrm{E}')\) 诱导拓扑的集合 B 到空间 \(\mathrm{F}_{\sigma}\) 的连续映射。而集合 \(u(\mathrm{B})\) 包含于 F 的一个紧子集 C（III, p. 2, 注记 1），且空间 \(\mathrm{F}_{\sigma}\) 是豪斯多夫的，故由 F 与由 \(\mathrm{F}_{\sigma}\) 在 C 上诱导的拓扑重合。命题得证。

推论. --- 设 \((x_{n})_{n\in\mathbf{N}}\) 是 E 的点的一个序列，它对弱拓扑收敛于 E 的一个点 x。则序列 \((u(x_{n}))_{n\in\mathbf{N}}\) 在 F 中收敛于 \(u(x)\)。

事实上，由点 \(x\) 与序列 \((x_{n})_{n\in \mathbf{N}}\) 的诸点所组成的集是 E 的一个有界子集（EVT, III, p. 3, 命题 2 之推论）。

命题 7. --- 设 E 是半赋范空间，E' 是它的对偶，B 是 E' 的单位球，u 是从 E' 到一个分离局部凸空间 F 的线性映射。若 u 在赋有 \(\sigma(\mathrm{E}',\mathrm{E})\) 诱导拓扑的 B 上的限制是连续的，则 \(u(\mathrm{B})\) 是 F 的一个紧子集，且 u 是从 E 的强对偶 \(\mathrm{E}_b'\) 到 F 的紧线性映射。

事实上，赋有 \(\sigma(\mathrm{E}^{\prime},\mathrm{E})\) 诱导拓扑的集合 B 是紧的（EVT, III, p. 17, cor. 2）。

推论 1. --- 设 E 是可数型的半赋范空间，E' 是它的对偶，B 是 E' 的单位球。设 u 是从 E' 到一个分离局部凸空间 F 的线性映射。假设对 B 中元素的每个序列 \((x_{n})_{n\in\mathbf{N}}\)（它对 \(\sigma(\mathrm{E}',\mathrm{E})\) 收敛于 0），序列 \((u(x_{n}))_{n\in\mathbf{N}}\) 都在 F 中收敛于 0。那么，\(u(\mathrm{B})\) 是 F 的一个紧子集，且 u 是从 E 的强对偶 \(E_{b}'\) 到 F 的紧线性映射。

事实上，\(\sigma(\mathrm{E}^{\prime},\mathrm{E})\) 在 B 上诱导的拓扑是可度量化的（EVT, III, p. 19, cor. 2）。因此，推论的假设蕴含：\(u\) 在 B 上的限制（当 B 赋有拓扑 \(\sigma(\mathrm{E}^{\prime},\mathrm{E})\) 时）是从 B 到 F 的连续映射，于是推论由前一命题得出。

推论 2. --- 设 E 是半赋范空间，S 是 E 的预紧子集所成的一个集，其并在 E 中稠密。用 \(E'_{b}\) 表示 E 的强对偶，用 \(E'_{S}\) 表示赋有 S-拓扑的 E 的对偶。空间 \(E'_{S}\) 是分离的，且从 \(E'_{b}\) 到 \(E'_{S}\) 的恒等映射是紧的。

依照 TG, X, p. 8, prop. 7，空间 \(E_{S}^{\prime}\) 是分离的；在 \(E^{\prime}\) 的单位球上，S-拓扑与弱拓扑 \(\sigma(\mathrm{E}^{\prime},\mathrm{E})\) 重合（EVT, III, p. 17, prop. 5），由此即得推论。

例. --- 设 E 是半赋范空间。从 \(E_{b}^{\prime}\) 到赋有弱拓扑、紧收敛拓扑或预紧收敛拓扑的空间 \(E^{\prime}\) 的恒等映射是紧的。

命题 8. --- 设 E 是自反的巴拿赫空间，B 是它的单位球，F 是分离的局部凸空间，u: E → F 是线性映射。下列条件等价：

\begin{enumerate}
\def\labelenumi{(\roman{enumi})}
\item
  线性映射 u 是紧的；
\item
  集合 \(u(\mathrm{B})\) 是 F 的一个紧子集；
\item
  u 在 B 上的限制是从赋有 \(\sigma(\mathrm{E},\mathrm{E}')\) 诱导拓扑的集合 B 到 F 的连续映射；
\item
  对 B 中点的任何序列 \((x_{n})_{n\in \mathbf{N}}\)（它对拓扑 \(\sigma (\mathrm{E},\mathrm{E}^{\prime})\) 收敛于 0），序列 \((u(x_n))_{n\in \mathbf{N}}\) 在 F 中收敛于 0；
\item
  从 \(u(\mathrm{B})\) 的点的任何无穷序列中都能取出一个在 F 中收敛的子序列；
\item
  \(u(\mathrm{B})\) 的点的任何无穷序列在 F 中都有聚点。
\end{enumerate}

由于 E 是自反空间，B 是 \(E_{\sigma}\) 的一个紧子集（EVT, IV, p. 17, prop. 6），故 (iii) 蕴含 (ii)。按定义，条件 (ii) 蕴含 (i)，而由命题 6，条件 (i) 蕴含 (iii)。同样初等的是 (iii) 蕴含 (iv) 且 (v) 蕴含 (vi)。

我们现在证明 (iv) 蕴含 (v)。设 \((x_{n})\) 是 B 的点的一个无穷序列。由于 B 是 \(E_{\sigma}\) 的紧子集，由什穆利安定理（EVT, IV, p. 36, th. 2），存在一个序列 \((y_{n})\)（从序列 \((x_{n})\) 中取出），它在 \(E_{\sigma}\) 中收敛于一个极限 y。序列 \((y_{n}-y)\) 在 \(E_{\sigma}\) 中有界，从而在 E 中有界（EVT, IV, p. 1, prop. 1）；因此它包含于 B 的一个位似集合。若条件 (iv) 成立，则序列 \((u(y_{n}-y))\) 在 F 中收敛于 0，而 \((u(y_{n}))\) 作为 \((u(x_{n}))\) 的一个子序列收敛于 \(u(y)\)。这就证明了 (iv) 蕴含 (v)。

最后，我们证明 (vi) 蕴含 (ii)。设 \(j: \mathrm{F} \to \widehat{\mathrm{F}}\) 是 F 到其完备化中的典范单射。在假设 (vi) 之下，\(u(\mathrm{B})\) 是 F 的一个预紧子集（EVT, IV, p. 32, prop. 1），因此 \(j(u(\mathrm{B}))\) 是 \(\widehat{\mathrm{F}}\) 的一个相对紧子集（TG, II, 第 2 目, p. 29），且 \(j \circ u\) 是紧线性映射。由已经证明的条件 (i) 与 (ii) 的等价性，\(j(u(\mathrm{B}))\) 是 \(\widehat{\mathrm{F}}\) 的一个紧子集。因此，\(u(\mathrm{B})\) 是 F 的一个紧子集。

\subsection{转置}

设 E 是局部凸空间，\(E^{\prime}\) 是它的对偶。回顾，我们用 \(E_{b}^{\prime}\) 和 \(E_{c}^{\prime}\) 分别表示给向量空间 \(E^{\prime}\) 赋以有界收敛拓扑和紧收敛拓扑所得到的局部凸空间（EVT, III, p. 14）。空间 \(E_{b}^{\prime}\) 也称为 E 的强对偶（同前引处, 例 4）。

命题 9. --- 设 E 是局部凸空间，F 是分离的局部凸空间，u 是从 E 到 F 的连续线性映射。

\begin{enumerate}
\def\labelenumi{\alph{enumi})}
\item
  若映射 \(u\) 是紧的，则它的转置 \(^{t}u\) 是从 \(F_c'\) 到 \(E_c'\) 的紧线性映射，且是从 \(F_c'\) 到 \(E_b'\) 的连续线性映射；
\item
  假设空间 E 是半赋范空间且空间 F 是拟完备的。若转置 \({}^{t}u\) 是从 \(F_{c}^{\prime}\) 到 \(E_{b}^{\prime}\) 的连续映射，则线性映射 u 是紧的。
\end{enumerate}

首先假设线性映射 \(u\) 是紧的。于是存在 E 中 0 的一个邻域 V，它在 \(u\) 下的像包含于 F 的一个紧子集 C。按照定义（EVT, II, p. 47, def. 2 及 p. 68, def. 1），C 的极集 \(\mathrm{C}^{\circ}\) 是 \(\mathrm{F}_c^\prime\) 中 0 的一个邻域，并且 \({}^t u(\mathrm{C}^\circ) \subset \mathrm{V}^\circ\)。而 \(\mathrm{V}^\circ\) 是 \(\mathrm{E}'\) 的一个等度连续且对弱拓扑为闭的子集，从而在 \(\mathrm{E}_c^\prime\) 中紧（EVT, III, p. 17, cor. 2 及 prop. 5）。这就证明了 \({}^t u\) 是从 \(\mathrm{F}_c^\prime\) 到 \(\mathrm{E}_c^\prime\) 的紧线性映射。

我们保留前面的假设，设 U 是 \(E_{b}^{\prime}\) 中 0 的一个邻域。它包含 E 的某个有界子集 A 的极集 \(A^{\circ}\)。由于线性映射 u 紧，集合 \(u(\mathrm{A})\) 在 F 中相对紧（III, p. 2, 注记 1），且 \((^{t}u)^{-1}(\mathrm{A}^{\circ}) = u(\mathrm{A})^{\circ}\) 是 \(F_{c}^{\prime}\) 中 0 的一个邻域。这就证明了 \(^{t}u\) 是从 \(F_{c}^{\prime}\) 到 \(E_{b}^{\prime}\) 的连续线性映射。

现在置于 \(b\) 的假设之下，并用 B 表示 E 的单位球。B 的极集 \(\mathrm{B}^{\circ}\) 是 \(\mathrm{E}_b'\) 的单位球，并且 \((^{t}u)^{-1}(\mathrm{B}^{\circ}) = u(\mathrm{B})^{\circ}\)。若 \(^{t}u\) 是从 \(\mathrm{F}_c'\) 到 \(\mathrm{E}_b'\) 的连续映射，则集合 \(u(\mathrm{B})^{\circ}\) 是 \(\mathrm{F}_c'\) 中 0 的一个邻域；于是它包含 F 的某个紧子集 C 的极集 \(\mathrm{C}^{\circ}\)。由双极定理（EVT, II, p. 49, cor. 3），集合 \(u(\mathrm{B})\) 包含于 \(\mathrm{C} \cup \{0\}\) 的闭凸包；由于空间 F 拟完备（EVT, III, p. 8），后者是紧的。这就证明了线性映射 \(u\) 是紧的。

推论 1（绍德尔）. --- 设 E 是半赋范空间，F 是巴拿赫空间，E' 和 F' 是它们各自的强对偶，u 是从 E 到 F 的连续线性映射。下列性质等价：

\begin{enumerate}
\def\labelenumi{(\roman{enumi})}
\item
  从 E 到 F 的线性映射 u 是紧的；
\item
  线性映射 \({}^{t}u\)（从 \(\mathrm{F}'\) 到 \(\mathrm{E}'\) 中）是紧的；
\item
  线性映射 \({}^{t}u\)（从 \(\mathrm{F}_{c}^{\prime}\) 到 \(\mathrm{E}'\) 中）是连续的。
\end{enumerate}

(i) 与 (iii) 的等价性由命题 9 的 a) 和 b) 得出；蕴含关系 (iii)⇒(ii) 来自从 F' 到 F'c 的恒等映射是紧的这一事实（III, p. 7, 命题 7 之推论）。

最后我们来证明条件 (ii) 蕴含 (i)。假设线性映射 \({}^{t}u\)：\(F^{\prime} \to E^{\prime}\) 是紧的。把蕴含关系 (i)⇒(ii) 应用于 \({}^{t}u\)，便知 \({}^{tt}u\) 是从 \(E''\) 到 \(F''\) 的紧线性映射。用 \(v\) 表示从 E 到 \(E''\) 的典范线性映射；它是连续的。由于 F 等同于 \(F''\) 的一个闭向量子空间，且 u 在 E 上与 \({}^{t}({}^{t}u) \circ v\) 重合，故由 III, p. 3, 注记 5 推出线性映射 \(u\) 是紧的。

推论 2. --- 设 E 是半赋范空间，F 是可数型的巴拿赫空间。设 u 是从 E 到 F 的连续线性映射。假设对 F' 中元素的每个弱趋于 0 的序列 \((y_{n})\)，序列 \((^{t}u(y_{n}))\) 都在 E' 中强趋于 0。则线性映射 u 是紧的。

用 B′ 表示 F′ 的单位球，赋有由拓扑 \(\sigma(\mathrm{F}',\mathrm{F})\) 诱导的拓扑。由于 F 是可数型的巴拿赫空间，B' 是紧的可度量化空间（EVT, III, p. 19, cor. 2）。在所作假设之下，\({}^{t}u\) 在 B' 上的限制是从 B' 到 E' 的连续映射，且 \({}^{t}u(\mathrm{B}')\) 是 E' 的一个紧子集。由推论 1，线性映射 u 是紧的。

注记. --- 设 E 和 F 是分离的局部凸空间。从 E 到 F 的紧线性映射的转置并不总是从 \(F_{b}^{\prime}\) 到 \(E_{b}^{\prime}\) 的紧线性映射（参见 III, p. 108, 习题 15）。

\subsection{勒雷--绍德尔定理}

下面的定理将在 TA（待出版）中证明。

定理 1（布劳威尔）. --- 设 B 是有限维赋范向量空间的一个非空紧凸子集。则从 B 到 B 的每个连续映射都有不动点。

我们将由它导出下面的结果。

定理 2（勒雷--绍德尔）. --- 设 X 是非空的紧拓扑空间，同胚于一个局部凸空间的凸子集。则从 X 到 X 的每个连续映射都有不动点。

我们可以假设 X 是局部凸空间 E 的一个凸子集。设 N 是 \(\{0\}\) 在 E 中的闭包，\(\pi: E \rightarrow E/N\) 是典范满射。由于 X 是分离空间，\(\pi\) 在 X 上的限制是单射；又由于 X 紧，它便定义一个从 X 到分离局部凸空间 E/N 的一个非空紧凸子集上的同胚。这就使我们能够假设空间 E 是分离的。

设 \(f: \mathrm{X} \to \mathrm{X}\) 是连续映射。映射 \(h: \mathrm{X} \to \mathrm{E}\)（由 \(h(x) = f(x) - x\) 定义）是连续的；由于 X 紧，像 \(h(\mathrm{X})\) 是闭的。因此只需证明 0 附着于 \(h(\mathrm{X})\)，也就是说，对 E 中 0 的每个开凸邻域 U，存在 X 的一个点 \(b\)，使得 \(f(b) - b \in \mathrm{U}\)。

设 U 是 E 中 0 的一个凸邻域。对每个 \(a \in X\)，用 \(V_a\) 表示 X 中这样的元素 \(x\) 所成的集：\(f(x) - f(a) \in U\)。它在 X 中是开的且包含 \(a\)。由于 X 紧且非空，存在 X 的一个非空有限子集 A，使得诸集合 \(V_a\)（\(a \in A\)）覆盖 X。设 \((\varphi_a)_{a \in A}\) 是从属于覆盖 \((V_a)_{a \in A}\) 的一个连续单位分解（TG, IX, p. 43, prop. 1 及 p. 47, th. 3）。用 F 表示由 \(f(A)\) 生成的 E 的子空间。对每个 \(x \in F \cap X\)，令

\[
g (x) = \sum_ {a \in A} \varphi_ {a} (x) f (a).
\]

这样就定义了一个连续映射 \(g: F \cap X \to E\)。它的像包含于 \(f(A)\) 的凸包，从而包含于 \(F \cap X\)。由于集合 \(F \cap X\) 是有限维向量空间的一个非空紧凸子集，由定理 1，映射 g 有一个不动点 \(b \in F \cap X\)。我们有

\[
f (b) - b = f (b) - g (b) = \sum_ {a \in \mathrm{A}} \varphi_ {a} (b) (f (b) - f (a)).
\]

对每个 \(a \in \mathrm{A}\)，若 \(\varphi_{a}(b) \neq 0\)，则我们有 \(b \in \mathrm{V}_{a}\)，从而 \(f(b) - f(a) \in \mathrm{U}\)。由于 U 是凸的，我们得到 \(f(b) - b \in \mathrm{U}\)，证明完毕。

\subsection{紧算子的不变子空间}

应当把下面的定理与舒尔引理（A, VIII, p. 43, cor. 及 V, p. 386, prop. 6）以及 I, p. 26 的推论 4 相比较。

定理 3. --- 设 E 是 C 上的分离局部凸空间，A 是 \(\mathcal{L}(\mathrm{E})\) 的一个子集。我们作下列假设：

\begin{enumerate}
\def\labelenumi{(\roman{enumi})}
\item
  不存在 E 的既异于 \(\{0\}\) 又异于 E 的、在 A 下稳定的闭向量子空间；
\item
  集合 A 含有一个非零的紧自同态。那么 A 的交换子化归于恒等映射的纯量倍数。
\end{enumerate}

用 A 生成的 \(\mathcal{L}(\mathrm{E})\) 的含单位元子代数代替 A，就化归为 A 是 \(\mathcal{L}(\mathrm{E})\) 的一个含单位元子代数的情形。于是设 \(h\) 是属于 A 的一个非零紧自同态。

引理 1. --- 存在 A 的一个元素 a，使得 \(ha - 1_{E}\) 的核不化归于 0。

存在 E 的一个元素 \(x_{0}\)，使得 \(h(x_{0}) \neq 0\)。设 V 是 \(h(x_{0})\) 的一个不包含 0 的闭邻域。由于自同态 h 紧，可以取 \(x_{0}\) 的一个开凸邻域 U，使得 \(h(\mathrm{U})\) 是 V 的一个相对紧子集。于是 \(h(\mathrm{U})\) 的闭包 C 是 E 的一个紧凸子集；由于它包含于 V，它不包含 0。

设 \(y\) 是 C 的一个点。用 \(Ay\) 表示 \(y\) 在 A 的诸元素下的像所成的集。它是 E 的一个在 A 下稳定的向量子空间；由于它包含 \(y\)，它是非零的。由定理的假设 (i)，集合 \(Ay\) 在 E 中稠密。因此存在 A 的一个元素 \(b\)，使得 \(b(y) \in U\)。由于集合 C 紧，存在 A 中元素的一个有限族 \((b_1, \ldots, b_n)\)，使得诸集合 \(b_i^{-1}(U)\) 覆盖 C。存在 C 上从属于这个覆盖的一个连续单位分解 \((\varphi_1, \ldots, \varphi_n)\)（TG, IX, p. 47, th. 3）。我们用下式定义映射 \(f: C \to E\)：

\[
f (x) = \sum_ {i = 1} ^ {n} \varphi_ {i} (x) b _ {i} (x)
\]

其中 \(x \in \mathrm{C}\)；这是一个连续映射。由于 U 是凸的，且 \(b_i(x)\) 在 \(\varphi_i(x)\) 非零时属于 U，我们有 \(f(\mathrm{C}) \subset \mathrm{U}\) 及 \(h(f(\mathrm{C})) \subset \mathrm{C}\)。由勒雷--绍德尔定理（III, p. 11, 定理 2），存在一个元素 \(x \in \mathrm{C}\)，使得 \(h(f(x)) = x\)。于是令

\[
a = \sum_ {i = 1} ^ {n} \varphi_ {i} (x) b _ {i}.
\]

我们有 \(h(a(x)) = h(f(x)) = x\)，且由于 0 不属于 C，x 非零，因此 \(ha - 1_{E}\) 的核非零。

我们来完成定理 3 的证明。设 \(a\) 是 A 的一个元素，使得 \(ha - 1_{\mathrm{E}}\) 的核 T 非零（引理 1）。用 \(i\) 表示 T 到 E 中的典范单射。线性映射 \(i\) 等于 \(h \circ a \circ i\)，因此它是紧的。于是 T 的维数有限（III, p. 2, 注记 3）。设 \(u\) 是 \(\mathcal{L}(\mathrm{E})\) 的一个与 A 交换的元素。由于 \(u\) 与 \(h \circ a\) 交换，我们有 \(u(\mathrm{T}) \subset \mathrm{T}\)。由于 T 在代数闭域 \(\mathbf{C}\) 上是有限维且非零的，\(u\) 在 T 上诱导的自同态有一个特征值 \(\lambda\)。

令 \(\mathrm{F} = \mathrm{Ker}(u - \lambda1_{\mathrm{E}})\)。这是 E 的一个非零闭向量子空间；由于 u 与 A 的诸元素交换，它在 A 下稳定。由假设 (i)，我们有 F = E，从而 \(u = \lambda1_{E}\)。

推论 1. --- 保持定理 3 的假设，并再设 A 的诸元素两两交换。那么 E 在 C 上是一维的。

由定理 3 的假设 (ii)，空间 E 非零。因此它包含一个维数为 1 的向量子空间 F。由于 E 假设为分离的，这个子空间是闭的。由定理 3，A 的每个元素都是位似变换，故使 F 稳定。于是由定理 3 的假设 (i) 推出 F 等于 E。

推论 2（洛莫诺索夫）. --- 设 E 是域 C 上维数至少为 2 的分离局部凸空间，u 是 E 的一个自同态。假设存在 E 的一个与 u 交换的紧自同态 \(h \neq 0\)。那么存在 E 的一个既异于 \(\{0\}\) 又异于 E 的闭向量子空间 F，使得 \(u(F) \subset F\)。

事实上，推论 1 表明集合 A = \{u, h\} 不可能满足定理 3 的假设 (i)。

推论 3. --- 设 E 是域 C 上的分离局部凸空间，u 是 E 的紧自同态。若 E 的维数至少为 2，则存在 E 的一个既异于 \(\{0\}\) 又异于 E 的闭向量子空间 F，使得 \(u(F) \subset F\)。

u = 0 的情形是直接的。\(u \neq 0\) 的情形由推论 2 取 h = u 得出。

\subsection{逼近空间}

给定局部凸空间 E 和 F，回顾（EVT, III, p. 14, 例 2），我们用 \(\mathcal{L}_{\mathrm{pc}}(\mathrm{E};\mathrm{F})\) 表示给 \(\mathcal{L}(\mathrm{E};\mathrm{F})\) 赋以预紧收敛拓扑所得到的局部凸空间。我们也用 \(\mathcal{L}_{\mathrm{pc}}(\mathrm{E})\) 表示局部凸空间 \(\mathcal{L}_{\mathrm{pc}}(\mathrm{E};\mathrm{E})\)。当局部凸空间 E 是豪斯多夫的且拟完备（EVT, III, p. 8, def. 6）时，\(\mathcal{L}(\mathrm{E};\mathrm{F})\) 上的预紧收敛拓扑与紧收敛拓扑重合，因为 E 的预紧子集的闭包是紧的（EVT, III, p. 8）。

定义 2. --- 局部凸空间 E 称为具有逼近性质，或又称为逼近空间，如果 E 的恒等映射 \(1_{E}\) 在空间 \(\mathcal{L}_{\mathrm{pc}}(\mathrm{E})\) 中附着于 E 的连续有限秩自同态所成的集 \(\mathcal{L}^{\mathrm{f}}(\mathrm{E})\)。

特别地，每个有限维局部凸空间都是逼近空间。

注记. --- 1) 为使局部凸空间 E 和 F 的拓扑直和是逼近空间，必须且只需 E 和 F 都是逼近空间。

\begin{enumerate}
\def\labelenumi{\arabic{enumi})}
\setcounter{enumi}{1}
\item
  设 E 是局部凸空间。设 N 是 \(\{0\}\) 在 E 中的闭包，P 是 N 在 E 中的一个代数补子空间。那么 P 是 N 在 E 中的拓扑补，且同构于局部凸空间 E/N。空间 N 是逼近空间。由注记 1，为使 E 是逼近空间，必须且只需豪斯多夫局部凸空间 E/N 是逼近空间。
\item
  设 E 是局部凸空间，A 是 \(\mathcal{L}^{\mathrm{f}}(\mathrm{E})\) 的一个等度连续子集，T 是 E 的一个完全子集。若 \(1_{\mathrm{E}}\) 对 T 上的逐点收敛拓扑附着于 A，则 \(1_{\mathrm{E}}\) 在 \(\mathcal{L}_{\mathrm{pc}}(\mathrm{E})\) 中附着于 A（EVT, III, p. 16, prop. 4 及 p. 17, prop. 5），从而 E 是逼近空间。
\item
  设 E 是 C 上的局部凸空间。用 \(E_{0}\) 表示 E 的底层 R 上的局部凸空间。为使 E 是逼近空间，必须且只需 \(E_{0}\) 本身是逼近空间。事实上，该条件是必要的；我们来证明它是充分的。于是设 \(E_{0}\) 是逼近空间。设 C 是 E 的一个预紧子集，U 是 E 中 0 的一个凸均衡邻域。令 \(C' = C \cup iC\)。存在一个从 \(E_{0}\) 到 \(E_{0}\) 的连续有限秩 R-线性映射 u，使得 \(x - u(x)\) 对一切 \(x \in C'\) 都属于 U。对 E 中的一切 x，令 \(v(x) = \frac{1}{2}(u(x) - iu(ix))\)。这样就定义了一个从 E 到 E 的连续有限秩 C-线性映射。对每个 \(x \in C\)，我们有 \(x \in C'\)，\(ix \in C'\)，并且 \(x - v(x) = \frac{1}{2}(x - u(x)) - \frac{i}{2}(ix - u(ix))\)，从而 \(x - v(x)\) 属于 U。这就证明了 E 是逼近空间。
\item
  设 E 是 R 上的局部凸空间。为使复化局部凸空间 \(\mathrm{E}_{(\mathbf{C})}\) 是逼近空间，必须且只需 E 本身是逼近空间。这由注记 1 和注记 4 得出，因为 \(\mathrm{E}_{(\mathbf{C})}\) 的底层实局部凸空间同构于 \(E \times E\)。
\end{enumerate}

命题 10. --- 设 E 是希尔伯特空间。那么 E 是逼近空间。

E 上的有限秩正交投影算子所成的集 A 是等度连续的。设 \(n \geqslant 1\) 是整数，\((x_{1}, \ldots, x_{n})\) 是 E 的元素的一个族。用 V 表示由诸 \(x_{i}\) 生成的向量子空间，用 \(p_{V}\) 表示以 V 为像的正交投影算子。我们有 \(p_{\mathrm{V}}(x_{i}) = x_{i}\)（对 \(1 \leqslant i \leqslant n\)）。这就蕴含 \(1_{E}\) 对逐点收敛拓扑附着于 A，从而 E 是逼近空间（注记 3）。

逼近空间的其他例子将在 III, p. 20, 第 7 目中给出。

引理 2. --- 设 E 是局部凸空间。那么 E 是逼近空间，当且仅当：对 E 的每个预紧子集 C 和 E 中 0 的每个邻域 U，存在一个整数 \(n \geqslant 0\)、E 的元素 \(e_{1}, \ldots, e_{n}\) 以及 E 上的连续线性形式 \(f_{1}, \ldots, f_{n}\)，使得 \(x - \sum_{i=1}^{n} f_{i}(x)e_{i}\) 对一切 \(x \in C\) 都属于 U。

预紧收敛拓扑的定义表明该条件是充分的。反之，设 E 是逼近空间。设 P 是 \(\{0\}\) 的闭包在 E 中的一个拓扑补空间（参见注记 2）。此时，\(u \in \mathcal{L}^{\mathrm{f}}(\mathrm{E})\)（满足 \(u(\mathrm{E}) \subset \mathrm{P}\)）所成的集 A 在 \(\mathcal{L}^{\mathrm{f}}(\mathrm{E})\) 中稠密。由于 P 是分离的，A 的每个元素 \(u\) 都形如 \(x \mapsto \sum_{i=1}^{n} f_i(x)e_i\)，其中 \(n \in \mathbf{N}\)，\(e_1, \ldots, e_n\) 是 E 的元素，而 \(f_1, \ldots, f_n\) 是 E 上的连续线性形式（EVT, I, p. 14, th. 2）。这就证明了该条件是必要的。

注记 6. --- 存在不具有逼近性质的巴拿赫空间，这是由 P. ENFLO 证明的（A counterexample to the approximation problem in Banach spaces, Acta Math. 130 (1973), p. 309--317; 参见 III, p. 112, 习题 25）。这回答了 S. BANACH 的一个问题（Théorie des opérations linéaires, Monografie Matematyczne I, Warszawa, 1932, 第 VI 章 § 1 中的注记）。另见 p. 21 上的注记。

命题 11. --- 设 E 和 F 是局部凸空间。假设 E 是逼近空间。

\begin{enumerate}
\def\labelenumi{\alph{enumi})}
\item
  集 \(\mathcal{L}^{\mathrm{f}}(\mathrm{E};\mathrm{F})\) 在 \(\mathcal{L}_{\mathrm{pc}}(\mathrm{E};\mathrm{F})\) 中稠密；
\item
  集 \(\mathcal{L}^{\mathrm{f}}(\mathrm{F};\mathrm{E})\) 在 \(\mathcal{L}_{\mathrm{pc}}(\mathrm{F};\mathrm{E})\) 中稠密；
\item
  设 \(\mathfrak{S}\) 是 F 的有界子集所成的一个集，\(u \in \mathcal{L}(\mathrm{F};\mathrm{E})\)。假设在 \(u\) 下，属于 \(\mathfrak{S}\) 的 F 的每个子集的像都是预紧的。那么 \(u\) 附着于 \(\mathcal{L}^{\mathrm{f}}(\mathrm{F};\mathrm{E})\)（对 \(\mathfrak{S}\)-拓扑）（EVT, III, p. 13）。
\end{enumerate}

设 \(v\) 是 \(\mathcal{L}(\mathrm{E};\mathrm{F})\) 的一个元素。映射 \(\varphi\colon w\mapsto v\circ w\)（从 \(\mathcal{L}_{\mathrm{pc}}(\mathrm{E})\) 到 \(\mathcal{L}_{\mathrm{pc}}(\mathrm{E};\mathrm{F})\) 中）是连续的（TG, X, p. 5, prop. 3）。我们有 \(\varphi(1_{\mathrm{E}})=v\) 及 \(\varphi(\mathcal{L}^{\mathrm{f}}(\mathrm{E}))\subset\mathcal{L}^{\mathrm{f}}(\mathrm{E};\mathrm{F})\)，因此 \(v\) 附着于 \(\mathcal{L}^{\mathrm{f}}(\mathrm{E};\mathrm{F})\)（在 \(\mathcal{L}_{\mathrm{pc}}(\mathrm{E};\mathrm{F})\) 中）。这就证明了 \(a\)）。

同样，在 \(c\)）的假设之下，映射 \(\psi: w \mapsto w \circ u\)（从 \(\mathcal{L}_{\mathrm{pc}}(\mathrm{E})\) 到 \(\mathcal{L}_{\mathfrak{S}}(\mathrm{F};\mathrm{E})\) 中）是连续的（同前引处）。我们有 \(\psi(1_{\mathrm{E}}) = u\) 及 \(\psi(\mathcal{L}^{\mathrm{f}}(\mathrm{E})) \subset \mathcal{L}^{\mathrm{f}}(\mathrm{F};\mathrm{E})\)，因此 \(u\) 附着于 \(\mathcal{L}^{\mathrm{f}}(\mathrm{F};\mathrm{E})\)（在 \(\mathcal{L}_{\mathfrak{S}}(\mathrm{F};\mathrm{E})\) 中）。

F 的预紧子集在从 F 到 E 的连续线性映射下的像是预紧的，故 b) 由 c) 得出。

推论. --- 设 E 是分离的逼近空间，F 是局部凸空间。从 F 到 E 的每个紧线性映射都附着于 \(\mathcal{L}^{\mathrm{f}}(\mathrm{F};\mathrm{E})\)（在赋有有界收敛拓扑的空间 \(\mathcal{L}(\mathrm{F};\mathrm{E})\) 中）。

这由命题 11 的 c) 得出，因为 F 的有界子集在从 F 到 E 的紧线性映射下的像在 E 中相对紧（III, p. 2, 注记 1），从而预紧。

命题 12. --- 设 E 是局部凸空间，I 是一个集合，且对每个 \(i \in I\)，设 \(F_{i}\) 是局部凸空间，\(v_{i}: E \to F_{i}\) 是有稠密像的连续线性映射。假设对 E 中 0 的每个邻域 U，存在 \(i \in I\) 及 \(F_{i}\) 中 0 的一个邻域 V，使得 \(v_{i}^{-1}(V) \subset U\)。若诸 \(F_{i}\) 都是逼近空间，则 E 也是逼近空间。

设 A 是 E 的一个预紧子集，U 是 E 中 0 的一个邻域。按假设，存在 \(i \in I\) 及一个连续半范数 \(p\)（定义在 \(F_i\) 上），使得 U 包含 \((p \circ v_i)^{-1}([0,1])\)。令 \(F = F_i\)，\(v = v_i\)，\(B = v(A)\)，并设 F 是逼近空间。集合 B 在 F 中预紧。因此，由引理 2，存在一个整数 \(n \geqslant 1\)、F 的元素 \(y_1, \ldots, y_n\) 以及 F 上的连续线性形式 \(f_1, \ldots, f_n\)，使得对一切 \(y \in B\)，有

\[
p \left(y - \sum_ {j = 1} ^ {n} f _ {j} (y) y _ {j}\right) \leqslant \frac {1}{2}.
\]

由于 B 有界（EVT, III, p. 3, prop. 2），存在实数 M > 0，使得 \(|f_{j}(y)| \leqslant M\) 对一切满足 \(1 \leqslant j \leqslant n\) 的 j 及每个 \(y \in B\) 都成立。此外，由于 \(v(\mathrm{E})\) 在 F 中稠密，对每个满足 \(1 \leqslant j \leqslant n\) 的整数 j，存在 E 的一个元素 \(x_{j}\)，使得 \(p(y_{j} - v(x_{j})) \leqslant (2n\mathrm{M})^{-1}\)。线性映射

\[
u \colon x \longmapsto \sum_ {j = 1} ^ {n} f _ {j} (v (x)) x _ {j}
\]

属于 \(\mathcal{L}^{\mathrm{f}}(\mathrm{E})\)。对每个 \(x\in \mathrm{A}\)，我们有

\[
v (x - u (x)) = v (x) - \sum_ {j = 1} ^ {n} f _ {j} (v (x)) y _ {j} + \sum_ {j = 1} ^ {n} f _ {j} (v (x)) \left(y _ {j} - v \left(x _ {j}\right)\right),
\]

由此 \(p(v(x - u(x))) \leqslant \frac{1}{2} + \frac{nM}{2nM} = 1\)，从而 \(x - u(x) \in U\)。命题得证。

推论 1. --- 逼近空间的一个稠密向量子空间是逼近空间。

推论 2. --- 若局部凸空间 E 的分离完备化是逼近空间，则 E 是逼近空间。

推论 3. --- 一族逼近空间的乘积是逼近空间。

事实上，设 \((\mathrm{E}_{i})_{i\in\mathrm{I}}\) 是一族逼近空间。对 I 的每个有限子集 J，令 \(E_{J}=\prod_{i\in J}E_{i}\)，并用 \(v_{J}\) 表示从 \(E=\prod_{i\in I}E_{i}\) 到 \(E_{J}\) 中的典范映射。局部凸空间 \(E_{J}\) 是逼近空间（注记 1）。于是，把命题 12 应用于局部凸空间 E、局部凸空间族 \((\mathrm{E}_{\mathrm{J}})\) 以及连续线性映射 \(v_{J}:E\to E_{J}\)，即得该推论。

推论 4. --- 设 E 是局部凸空间。若 E 上的每个连续半范数都被一个连续预希尔伯特半范数所控制，则 E 是逼近空间。

设 \(\mathcal{P}\) 是 E 上连续预希尔伯特半范数所成的集。对每个 \(p \in \mathcal{P}\)，假设蕴含：半赋范空间 \(\mathrm{E}_p\)（给 E 赋以半范数 \(p\) 而得到）是逼近空间（推论 2 与命题 10），并且从 E 到 \(\mathrm{E}_p\) 的恒等映射是连续的。因此该推论由命题 12 得出。

引理 3. --- 设 E 是可度量化的局部凸空间，\((x_{n})_{n\in\mathbf{N}}\) 是 E 中收敛于 0 的元素的一个序列。则存在 E 中收敛于 0 的元素的一个序列 \((y_{n})_{n\in\mathbf{N}}\)，以及区间 [0,1] 中收敛于 0 的元素的一个序列 \((\lambda_{n})_{n\in\mathbf{N}}\)，使得 \(x_{n}=\lambda_{n}y_{n}\) 对一切 \(n\in N\) 成立。

由于 E 可度量化，存在 E 中 0 的均衡邻域的一个基本系 \((\mathrm{V}_{m})_{m\in\mathbf{N}}\)，使得 \(V_{0}=E\)，且 \(2^{m+1}V_{m+1}\subset V_{m}\) 对一切 \(m\geqslant0\) 成立。由于 \((x_{n})\) 收敛于 0，存在整数的一个严格递增序列 \((\mathrm{N}_{m})_{m\in\mathbf{N}}\)，使得 \(N_{0}=0\)，且对一切 \(m\geqslant0\)，有 \(x_{n}\in V_{m}\)（对 \(n\geqslant N_{m}\)）。对每个整数 \(n\geqslant0\)，存在唯一的整数 \(m\geqslant0\)，使得 \(N_{m}\leqslant n<N_{m+1}\)，这时我们令 \(\lambda_{n}=2^{-m}\)，\(y_{n}=2^{m}x_{n}\)。如此定义的序列 \((\lambda_{n})\) 收敛于 0。此外，由于 \(y_{n}\in V_{m}\)（对 \(n\geqslant N_{m+1}\)），序列 \((y_{n})\) 在 E 中收敛于 0。最后，对一切 n，我们有 \(x_{n}=\lambda_{n}y_{n}\)。

回顾（EVT, II, p. 28），若 E 是向量空间而 L 是 E 的一个凸均衡子集，我们用 \(E_{L}\) 表示由 L 生成的 E 的向量子空间，赋以以 L 为单位球的半范数。当集合 L 不含任何直线时，这个半范数是一个范数。

引理 4. — 设 E 是弗雷歇空间，A 是 E 的一个紧子集。则存在 E 的一个包含 A 的紧凸均衡子集 L，使得由 E 与由 \(E_{L}\) 在 A 上诱导的拓扑重合。

集合 A 包含于 E 中收敛于 0 的元素的一个序列 \((x_{n})_{n\in\mathbf{N}}\) 的诸点的闭凸均衡包中（EVT, IV, p. 24, cor. 1）。设 \((y_{n})_{n\in\mathbf{N}}\) 与 \((\lambda_{n})_{n\in\mathbf{N}}\) 是满足引理 3 的结论的序列。序列 \((y_{n})\) 的诸点的闭凸均衡包 L 包含 A，且是 E 的一个紧子集（EVT, III, p. 8）。因此，\(E_{L}\) 是巴拿赫空间（EVT, III, p. 8, cor.）。在这个巴拿赫空间中，\(x_{n}\) 的范数被 \(\lambda_{n}\) 所控制，故序列 \((x_{n})\) 趋于 0。\(\widetilde{A}\)（序列 \((x_{n})\) 在 \(E_{L}\) 中的闭凸均衡包）是 \(E_{L}\) 的一个紧子集（EVT, III, p. 8）。由于 L 是 E 的有界子集，\(E_{L}\) 到 E 中的典范单射是连续的。在 \(\widetilde{A}\) 上，由 \(E_{L}\) 诱导的拓扑与由 E 诱导的拓扑重合，且 \(\widetilde{A}\) 是 E 的一个紧子集，从而是闭子集。由于集合 \(\widetilde{A}\) 是凸的、均衡的且包含序列 \((x_{n})\)，故它包含 A。证明完毕。

定理 4. --- 设 E 是弗雷歇空间。

\begin{enumerate}
\def\labelenumi{\alph{enumi})}
\item
  假设 E 是逼近空间。那么，对每个半赋范空间 F，空间 \(\mathcal{L}^{\mathrm{f}}(\mathrm{F};\mathrm{E})\) 对有界收敛拓扑在 \(\mathcal{L}^{\mathrm{c}}(\mathrm{F};\mathrm{E})\) 中稠密；
\item
  反之，假设对每个巴拿赫空间 F，从 F 到 E 的每个紧线性映射都对有界收敛拓扑附着于 \(\mathcal{L}^{\mathrm{f}}(\mathrm{F};\mathrm{E})\)。那么 E 是逼近空间。
\end{enumerate}

设 F 是半赋范空间。\(\mathcal{L}^{\mathrm{f}}(\mathrm{F};\mathrm{E})\) 在 \(\mathcal{L}(\mathrm{F};\mathrm{E})\) 中的闭包含于 \(\mathcal{L}^{\mathrm{c}}(\mathrm{F};\mathrm{E})\)（III, p. 4, 命题 1 与命题 2）。若 E 是逼近空间，则由命题 11 的推论，它等于 \(\mathcal{L}^{\mathrm{c}}(\mathrm{F};\mathrm{E})\)。这就证明了断言 \(a\)）。

假设 \(b\)）的假设成立。设 \(\varepsilon > 0\) 是实数。设 A 是 E 的一个紧子集，\(p\) 是 E 上的一个连续半范数。设 L 是 E 的一个紧凸均衡子集，使得 A 是赋范空间 \(\mathrm{E_L}\) 的一个紧子集（引理 4）。典范单射 \(j: \mathrm{E_L} \to \mathrm{E}\) 是紧的，且 \(\mathrm{E_L}\) 是巴拿赫空间（EVT, III, p. 8, cor.）。因此，按假设，存在一个整数 \(n \geqslant 1\)、E 的元素 \(e_1, \ldots, e_n\) 以及连续线性形式 \(\ell_1, \ldots, \ell_n\)（定义在 \(\mathrm{E_L}\) 上），使得映射 \(v\)（它从 \(\mathrm{E_L}\) 到 E，由 \(v(x) = \sum_{i=1}^{n} \ell_i(x)e_i\) 定义）满足 \(p(x - v(x)) \leqslant \frac{\varepsilon}{2}\)（对 \(x \in A\)）。映射 \({}^{t}j: E' \to (E_{L})'\) 的像对弱拓扑在 \((E_{L})'\) 中稠密（EVT, IV, p. 6, prop. 5）。在 \((E_{L})'\) 上，紧收敛拓扑与 \((E_{L})'\) 和 \(E_{L}\) 之间的对偶相容（EVT, IV, p. 3, 例）。因此 \({}^{t}j(E')\) 对紧收敛拓扑在 \((E_{L})'\) 中稠密（EVT, IV, p. 1, prop. 1），从而存在 E 上的连续线性形式 \(f_{1}, \ldots, f_{n}\)，使得

\[
| \ell_ {i} (x) - f _ {i} (x) | p (e _ {i}) \leqslant \frac {\varepsilon}{2 n}
\]

其中 \(x \in \mathrm{A}\)，\(1 \leqslant i \leqslant n\)。自同态 \(u \colon x \mapsto \sum_{i=1}^{n} f_i(x)e_i\)（\(\mathrm{E}\) 的）属于 \(\mathcal{L}^{\mathrm{f}}(\mathrm{E})\)，且对每个 \(x \in \mathrm{A}\)，我们有

\[
p (x - u (x)) \leqslant p (x - v (x)) + p (v (x) - u (x)) \leqslant \frac {\varepsilon}{2} + n \times \frac {\varepsilon}{2 n} = \varepsilon .
\]

由此推出，\(1_{E}\) 对紧收敛拓扑附着于 \(\mathcal{L}^{\mathrm{f}}(\mathrm{E})\)。而由于 E 完备，这个拓扑与预紧收敛拓扑重合。因此 E 是逼近空间。

\subsection{逼近空间的例子}

回顾，每个希尔伯特空间都是逼近空间（III, p. 15, 命题 10）。本节给出其它一些例子。

若 X 是局部紧空间，我们用 \(\mathcal{C}_{0}(\mathrm{X};\mathrm{K})\)，或简记作 \(\mathcal{C}_{0}(\mathrm{X})\)，表示在无穷远处消失的从 X 到 K 的连续函数所构成的巴拿赫空间，赋以由 \(\|f\|=\sup_{x\in\mathrm{X}}|f(x)|\)（对 \(f\in\mathcal{C}_{0}(\mathrm{X})\)）定义的范数。当 X 紧时，这个空间与从 X 到 K 的连续函数构成的空间 \(\mathcal{C}(\mathrm{X})\) 重合。

引理 5. --- 设 X 是紧拓扑空间，F 是 \(\mathcal{C}(\mathrm{X})\) 的一个有限子集，\(\varepsilon > 0\) 是实数。

存在一个有限子集 \(X_{0} \subset X\) 和一个线性映射 \(u: \mathcal{C}(X_{0}) \to \mathcal{C}(X)\)（范数 \(\leqslant 1\)），使得 \(\|u(f|X_{0}) - f\| \leqslant \varepsilon\) 对一切 \(f \in F\) 成立。

由于集合 F 是 \(\mathcal{C}(\mathrm{X})\) 的等度连续子集，存在 X 的一个有限覆盖 \((\mathrm{U}_{i})_{i\in\mathrm{I}}\)，使得对每个 \(i\in I\) 和每个 \(f\in F\)，\(f(\mathrm{U}_{i})\) 的直径 \(\leqslant\varepsilon\)。对每个 \(i\in I\)，选取一个点 \(x_{i}\)（属于 \(U_{i}\)），并用 \(X_{0}\) 表示诸 \(x_{i}\)（\(i\in I\)）所成的有限集。

设 \((\varphi_{i})_{i\in\mathrm{I}}\) 是从属于 \((\mathrm{U}_{i})_{i\in\mathrm{I}}\) 的连续单位分解。如下定义一个从 \(\mathcal{C}(\mathrm{X}_{0})\) 到 \(\mathcal{C}(\mathrm{X})\) 中的线性映射 u：令

\[
u (g) = \sum_ {i \in \mathrm{I}} g (x _ {i}) \varphi_ {i}
\]

其中 \(g \in \mathcal{C}(X_{0})\)。线性映射 u 的范数 \(\leqslant 1\)，且对 F 中的一切 f 满足 \(\|u(f|X_{0}) - f\| \leqslant \varepsilon\)，由此即得断言。

命题 13. --- 设 X 是局部紧拓扑空间。巴拿赫空间 \(\mathcal{C}_{0}(\mathrm{X})\) 是逼近空间。

首先假设 X 紧。设 \(A \subset \mathcal{L}^{\mathrm{f}}(\mathcal{C}(\mathrm{X}))\) 是形如 \(f \mapsto u(f|X_{0})\) 的映射所成的集合，其中 \(X_{0}\) 是 X 的有限子集，而 \(u : \mathcal{C}(X_{0}) \to \mathcal{C}(X)\) 是范数 \(\leqslant 1\) 的线性映射。集合 A 等度连续。由引理 5，\(\mathcal{C}(X)\) 的恒等映射对 \(\mathcal{C}(X)\) 上的逐点收敛拓扑附着于 A。于是命题由 III, p. 14 的注记 3 得出。

转到一般情形。设 Y 是 X 的亚历山德罗夫紧化，\(\omega\) 是 Y 的无穷远点（TG, I, p. 67）。把 \(\mathcal{C}_{0}(X)\) 与 \(\mathcal{C}(Y)\) 中在 \(\omega\) 处消失的元素所成的集合等同起来。那么 \(\mathcal{C}(Y)\) 是 \(\mathcal{C}_{0}(X)\) 与 Y 上常值函数的向量空间的拓扑直和。由于按前面所证 \(\mathcal{C}(Y)\) 是逼近空间，故空间 \(\mathcal{C}_{0}(X)\) 是逼近空间（III, p. 14 的注记 1）。

推论. --- 每个交换星代数都是逼近空间。

事实上，这样的代数同构于一个局部紧空间上在无穷远处消失的连续函数所构成的代数（I, p. 108, 定理 1）。

注记. --- 若 E 是无穷维希尔伯特空间，则星代数 \(\mathcal{L}(\mathrm{E})\) 不具有逼近性质（A. SZANKOWSKI, \(\mathcal{B}(\mathrm{H})\) does not have the approximation property, Acta Math. 147 (1981), p. 89--108）。

引理 6. --- 设 X 是局部紧拓扑空间。设 \(\mu\) 是 X 上的一个正测度，\(p \in [1, +\infty[\)。设 F 是 \(\mathrm{L}_{\mathrm{K}}^{p}(\mathrm{X}, \mu)\) 的一个有限子集，\(\varepsilon > 0\) 是实数。

存在 \(\mathrm{L}_{\mathrm{K}}^{p}(\mathrm{X},\mu)\) 的一个范数 \(\leqslant 1\) 的有限秩投影子 u，使得对 F 中的每个 f，\(\|u(f)-f\|\leqslant\varepsilon\)。

设 \(\mathcal{P}\) 是 X 的有限分割 \(\pi = (\mathrm{K}_1, \ldots, \mathrm{K}_n, \mathrm{H})\) 所成的集合，其中 \(n \geqslant 1\) 是整数，而 \(\mathrm{K}_1, \ldots, \mathrm{K}_n\) 是 X 的测度非零的可积子集。对每个分割 \(\pi = (\mathrm{K}_1, \ldots, \mathrm{K}_n, \mathrm{H}) \in \mathcal{P}\)，如下定义自同态 \(v_\pi\)（\(\mathcal{L}^p(\mathrm{X}, \mu)\) 的）：令

\[
v _ {\pi} (f) = \sum_ {i = 1} ^ {n} \frac {1}{\mu (\mathrm{K} _ {i})} \left(\int_ {\mathrm{K} _ {i}} f d \mu\right) \varphi_ {\mathrm{K} _ {i}},
\]

其中 \(f \in \mathcal{L}_{\mathrm{K}}^{p}(\mathrm{X}, \mu)\)，\(\varphi_{\mathrm{K}_i}\) 是 \(\mathrm{K}_i\) 的特征函数。映射 \(v_\pi\) 通过过渡到商，诱导出一个投影子 \(u_\pi\)（\(\mathrm{L}_{\mathrm{K}}^{p}(\mathrm{X}, \mu)\) 的）。容易验证，\(u_\pi\) 的像是由满足以下条件的函数 \(f \in \mathcal{L}_{\mathrm{K}}^{p}(\mathrm{X}, \mu)\) 的类所构成的空间：\(f\) 在 H 上为零，且在 \(\mathrm{K}_i\) 上取常数（对 \(1 \leqslant i \leqslant n\)）。

我们来证明 \(\| u_{\pi}\| \leqslant 1\)。对 \(f\in \mathcal{L}_{\mathrm{K}}^{p}(\mathrm{X},\mu)\)，我们有

\[
\| u _ {\pi} (f) \| _ {p} ^ {p} = \sum_ {i = 1} ^ {n} \mu (\mathrm{K} _ {i}) ^ {1 - p} \left| \int_ {\mathrm{K} _ {i}} f d \mu \right| ^ {p}.
\]

由赫尔德不等式（INT, IV, p. 208, § 6, 第 4 目, 定理 2），对每个 i 我们有不等式

\[
\left| \int_ {\mathrm{K} _ {i}} f d \mu \right| ^ {p} \leqslant \mu (\mathrm{K} _ {i}) ^ {p - 1} \int_ {\mathrm{K} _ {i}} | f | ^ {p} d \mu ,
\]

由此得 \(\| u_{\pi}(f)\| _p\leqslant \| f\| _p.\)

设 \(\mathcal{E}\) 是 X 上只取有限多个值的可积函数在 \(\mathrm{L}_{\mathrm{K}}^{p}(\mathrm{X},\mu)\) 中的类所成的集合。由于 \(\mathcal{E}\) 在 \(\mathrm{L}_{\mathrm{K}}^{p}(\mathrm{X},\mu)\) 中稠密（INT, IV, p. 162, §4, 第 10 目, 推论 1），存在一个有限集 \(\mathrm{F}'\)（含于 \(\mathcal{E}\)），使得 F 的每个元素都以至多 \(\varepsilon\) 的距离接近 \(\mathrm{F}'\) 的某个元素。考虑由诸测度非零的集合所构成的有限分割 \(\pi\)，映射 \(x\mapsto (f(x))_{f\in \mathrm{F}'}\) 在每个这样的集合上取一个给定的值，便可看出存在一个元素 \(\pi\)（属于 \(\mathcal{P}\)），使得 \(u_{\pi}(f) = f\) 对一切 \(f\)（含于 \(\mathrm{F}'\)）成立。因此投影子 \(u_{\pi}\) 具有所要的性质。

命题 14. --- 设 X 是局部紧拓扑空间，\(\mu\) 是 X 上的一个正测度，\(p \in [1, +\infty]\)。空间 \(\mathrm{L}_{\mathrm{K}}^{p}(\mathrm{X}, \mu)\) 是逼近空间。

若 \(p = +\infty\)，则空间 \(\mathrm{L}_{\mathbf{C}}^{\infty}(\mathrm{X},\mu)\) 是交换星代数（I, p. 103, 例 4），因而是逼近空间（命题 13 的推论），而由 III, p. 15 的注记 5，\(\mathrm{L}_{\mathbf{R}}^{\infty}(\mathrm{X},\mu)\) 亦然。

假设 p 有限。设 \(\mathrm{A}\subset\mathcal{L}^{\mathrm{f}}(\mathrm{L}_{\mathrm{K}}^{p}(\mathrm{X},\mu))\) 是范数 \(\leqslant1\) 的有限秩投影子所成的集合。由引理 6，\(\mathrm{L}_{\mathrm{K}}^{p}(\mathrm{X},\mu)\) 的恒等映射对逐点收敛拓扑附着于 A，于是命题由 III, p. 14 的注记 3 得出。

\section{紧线性映射的例子}

\subsection{\texorpdfstring{有限迹、希尔伯特--施密特与核型 \(p^{e}\) 幂型的自同态}{有限迹、希尔伯特--施密特与核型 p\^{}\{e\} 幂型的自同态}}

*设 E 与 F 是希尔伯特空间。E 的每个有限迹自同态（EVT, V, p. 49, 定义 7 与 p. 50, 定义 8）都是紧的（IV, p. 165, 推论 1）；从 E 到 F 的每个希尔伯特--施密特映射（EVT, V, p. 51, 定义 9）都是紧的（IV, p. 165, 推论 2）。

对每个实数 \(p \geqslant 1\)，向量空间 \(\mathcal{L}_{p}(\mathrm{E};\mathrm{F})\)（由从 E 到 F 的核型 \(p^{\mathrm{e}}\) 幂型线性映射构成）包含于 \(\mathcal{L}^{\mathrm{c}}(\mathrm{E};\mathrm{F})\)（对 \(p = 1\) 的情形参看 IV, p. 169, 注记 2）。*

\subsection{序列空间中的对角算子}

设 I 是非空集合。回顾（EVT, I, p. 4 与 p. 5），由有界族 \(x = (x_i)_{i \in \mathrm{I}}\)（\(\mathbf{C}\) 的元素）构成的向量空间，赋以范数 \(\| x \| = \sup_i |x_i|\)，是记作 \(\ell_{\mathbf{C}}^{\infty}(\mathrm{I})\) 的巴拿赫空间。

设 \(p\) 是 \([1, +\infty[\) 的一个元素。空间 \(\ell_{\mathbf{C}}^{p}(\mathrm{I})\)（由族 \(x = (x_{i})_{i \in \mathrm{I}}\)（\(\mathbf{C}\) 的元素，且使族 \((|x_{i}|^{p})\) 可和）构成）是一个向量空间，在其上映射 \(x \mapsto \| x \| = (\sum_{i} |x_{i}|^{p})^{1/p}\) 是一个范数，且使它成为巴拿赫空间。若给集合 I 赋以离散拓扑以及使 \(\mu(\{x\}) = 1\)（对一切 \(x \in \mathrm{I}\)）的测度，这个空间正是巴拿赫空间 \(\mathrm{L}_{\mathbf{C}}^{p}(\mathrm{I}, \mu)\)。当 \(p = 1\) 或 p = 2 时，这个记号于是与 EVT, I, p. 4 及 EVT, V, p. 18 中的记号一致（参看 INT, IV, p. 141, §4, 第 1 目, 例）。

对复数的族 \(x = (x_{i})_{i \in \mathrm{I}}\) 与 \(y = (y_{i})_{i \in \mathrm{I}}\)，用 xy 表示族 \((x_{i}y_{i})_{i \in \mathrm{I}}\)。

命题 1. --- 设 \(\lambda = (\lambda_i)_{i\in I}\) 是 C 的元素的一个有界族。设 \(p\) 是 \([1, +\infty]\) 的一个元素。设 E 是巴拿赫空间 \(\ell_{\mathbf{C}}^p (\mathrm{I})\)。

\begin{enumerate}
\def\labelenumi{\alph{enumi})}
\item
  对每个 \(x \in E\)，有 \(\lambda x \in E\)，且映射 u：\(x \mapsto \lambda x\) 是 E 的连续自同态，它在巴拿赫代数 \(\mathcal{L}(E)\) 中的谱等于诸 \(\lambda_{i}\) 所成集合在 C 中的闭包，其范数等于 \(\sup_{i} |\lambda_{i}|\)；
\item
  自同态 u 是紧的，当且仅当对每个实数 \(\varepsilon > 0\)，元素 \(i \in I\) 中满足 \(|\lambda_{i}| \geqslant \varepsilon\) 者所成的集合是有限的。
\end{enumerate}

采用命题中的记号，自同态 \(x \mapsto \lambda x\) 称为 \(\ell_{\mathbf{C}}^{p}(\mathrm{I})\) 的对角自同态。

我们来证明该命题。设 \(C = \sup_{i} |\lambda_{i}|\)。设 \(x \in E\)。我们有不等式

\[
\sum_ {i \in I} | \lambda_ {i} x _ {i} | ^ {p} \leqslant C ^ {p} \| x \| ^ {p}, \text {   if   } p \neq + \infty
\]

\[
\sup _ {i \in I} | \lambda_ {i} x _ {i} | \leqslant C \| x \|, \text {   if   } p = + \infty .
\]

这就证明了 \(\lambda x \in E\)。于是映射 \(x \mapsto \lambda x\) 是 E 的自同态，且这些不等式证明它具有范数 \(\leqslant C\)。

对 \(j \in \mathrm{I}\)，用 \(e_j\) 表示元素 \((x_i)_{i \in \mathrm{I}}\)（\(\ell_{\mathbf{C}}^p(\mathrm{I})\) 中的）：它满足 \(x_j = 1\)，且 \(x_i = 0\)（当 \(i \neq j\)）。于是有 \(u(e_j) = \lambda_j e_j\)（对一切 \(j \in \mathrm{I}\)），这表明 \(\lambda_j\) 属于 \(u\) 的谱。由于 \(u\) 的谱是闭的，在 \(\mathbf{C}\) 中诸 \(\lambda_i\) 所成集合的闭包包含于 \(u\) 的谱。由于 \(u\) 的谱包含在以 0 为中心、以 \(\|u\|\) 为半径的圆盘中（I, p. 24, 定理 1 与 I, p. 21 的公式 (3)），我们推出不等式 \(C \leqslant \|u\|\)，从而 \(\|u\| = C\)。

反之，若 \(\lambda \in \mathbf{C}\) 不附着于诸 \(\lambda_{i}\) 所成的集合，则族 \(((\lambda_{i} - \lambda)^{-1})_{i\in \mathrm{I}}\) 有界。于是前面的论证表明，线性映射 \(x\mapsto ((\lambda -\lambda_i)^{-1}x_i)_{i\in \mathrm{I}}\) 是 E 的一个自同态。它是 \(u - \lambda 1_{\mathrm{E}}\) 的逆，因而 \(\lambda\) 不属于 \(u\) 的谱。这就证明了 \(a\)）。

现在证明 \(b\)）。假设自同态 \(u\) 是紧的。设 \(\varepsilon > 0\)。用 J 表示 \(i \in I\) 中满足 \(|\lambda_i| \geqslant \varepsilon\) 者所成的集合。

诸元素 \(e_i\)（\(i \in \mathbf{J}\)）所成的集合在 E 中有界。它在 \(u\) 下的像由诸元素 \(\lambda_i e_i\)（\(i \in \mathbf{J}\)）组成，因此在 E 中相对紧（III, p. 2 的注记 1）。由于 \(\| \lambda_i e_i - \lambda_j e_j \| \geqslant \varepsilon\)（对 J 中每一对元素 \(i \neq j\)），这仅当集合 J 有限时才可能。

我们来证明逆命题。设 J 是 I 的一个有限子集，记 \(z_{\mathrm{J}} = (z_{\mathrm{J},i})_{i \in \mathrm{I}}\)，其中 \(z_{J,i} = \lambda_i\)（当 \(i \in J\)），而 \(z_{J,i} = 0\)（当 \(i \in I - J\)）；族 \(z_J\) 定义了 E 的一个有限秩自同态 \(u_J \colon x \mapsto z_J x\)。由前所证，线性映射 \(u - u_J\) 的范数被 \(R_+\) 中族 \((|\lambda_i|)_{i \in I-J}\) 的上确界所控制。该假设蕴含：对每个 \(\varepsilon > 0\)，存在有限子集 \(J \subset I\)，使得这个上界 \(\leqslant \varepsilon\)。由此推出映射 u 附着于 \(\mathcal{L}^{\mathrm{f}}(\mathrm{E})\)，因而它是紧的（III, p. 4, 命题 2 的推论）。

注记. --- 1) 若 I 有限，断言 b) 的假设总成立。当集合 I 无穷时，该假设意味着族 \((\lambda_{i})_{i\in\mathrm{I}}\) 沿 I 的有限子集的余集所成的滤子趋于 0；特别地，\(i\in I\) 中满足 \(\lambda_{i}\neq0\) 者所成的集合这时是可数的。

\begin{enumerate}
\def\labelenumi{\arabic{enumi})}
\setcounter{enumi}{1}
\tightlist
\item
  我们稍后将看到（参看 IV, p. 149, 定理 1），当 \(p = 2\) 时，希尔伯特型空间 \(\ell_{\mathbf{C}}^{2}(\mathrm{I})\) 的每个紧正规自同态都形如 \(u = w \circ v \circ w^{-1}\)，其中 \(v\) 是紧对角自同态，而 \(w\) 是希尔伯特空间 \(\ell_{\mathbf{C}}^{2}(\mathrm{I}).*\) 的酉自同态（EVT, V, p. 40）。
\end{enumerate}

\subsection{由核定义的取值于连续函数空间中的线性映射}

这里我们取 K = C。设 X 与 Y 是局部紧拓扑空间。设 \(\mu\) 是 X 上的一个复测度，k 是从 X × Y 到 C 中的、具有下列两条性质的映射：

\begin{enumerate}
\def\labelenumi{(\roman{enumi})}
\item
  对每个 \(y \in Y\)，从 X 到 C 中的函数 \(k_{y} \colon x \mapsto k(x, y)\) 是 \(\mu\) 可积的；
\item
  映射 \(y \mapsto k_{y}\) 从 Y 到 \(\mathcal{L}_{\mathbf{C}}^{1}(\mathrm{X}, \mu)\) 中是连续的。
\end{enumerate}

给 Y 上的复值连续函数空间 \(\mathcal{C}(\mathrm{Y})\) 赋以紧收敛拓扑。对 \(f\)（\(\mathcal{L}_{\mathbf{C}}^{\infty}(\mathrm{X},\mu)\) 中的），定义 \(\widetilde{k}(f)\colon \mathrm{Y}\to \mathbf{C}\) 如下：

\[
\widetilde {k} (f) (y) = \int_ {\mathrm{X}} k (x, y) f (x) d \mu (x) = \int_ {\mathrm{X}} k _ {y} f d \mu .\tag{1}
\]

映射 \(h \mapsto \int_{X} hfd\mu\) 是 \(\mathcal{L}_{\mathbf{C}}^{1}(X,\mu)\) 上的连续线性形式，故条件 (ii) 表明函数 \(\widetilde{k}(f)\) 是连续的。

命题 2. — 映射 \(\widetilde{k}\)（从 \(\mathcal{L}_{\mathbf{C}}^{\infty}(\mathrm{X},\mu)\) 到 \(\mathcal{C}(\mathrm{Y})\) 中）是紧的。

对 Y 中任意的 \(y\) 与 \(y'\)，我们有

\[
| \widetilde {k} (f) (y) | \leqslant \| k _ {y} \| _ {1} \| f \| _ {\infty}, \quad | \widetilde {k} (f) (y) - \widetilde {k} (f) (y ^ {\prime}) | \leqslant \| k _ {y} - k _ {y ^ {\prime}} \| _ {1} \| f \| _ {\infty}.
\]

由阿斯科利定理（TG, X, p. 18, 推论 2）可得，在 \(\widetilde{k}\) 下，\(\mathcal{L}_{\mathbf{C}}^{\infty}(\mathrm{X};\mu)\) 的单位球的像在 \(\mathcal{C}(\mathrm{Y})\) 中相对紧。

推论. --- 设 X 是紧拓扑空间，Y 是局部紧拓扑空间，\(\mu\) 是 X 上的一个复测度，k 是从 X × Y 到 C 中的连续映射。那么公式 (1) 定义了一个仍记作 \(\widetilde{k}\) 的紧线性映射，它从赋以一致收敛拓扑的 \(\mathcal{C}(X)\) 到赋以紧收敛拓扑的 \(\mathcal{C}(Y)\) 中。

我们来验证 \(k\) 的条件 (i) 与 (ii)。对每个 \(y \in \mathrm{Y}\)，映射 \(k_y\) 在紧空间 X 上连续，因而是 \(\mu\) 可积的。此外，由于 X 紧（TG, X, p. 28, 定理 3），映射 \(y \mapsto k_y\) 从 Y 到 \(\mathcal{C}(\mathrm{X})\) 中连续。由于 \(\| k_y - k_{y'} \|_1 \leqslant \| \mu \| \| k_y - k_{y'} \|\)（对一切 \((y, y') \in \mathrm{Y}^2\)），这个从 Y 到 \(\mathcal{L}_{\mathbf{C}}^1(\mathrm{X}, \mu)\) 中的映射是连续的。由命题 2，映射 \(\widetilde{k}\)（从 \(\mathcal{L}_{\mathbf{C}}^\infty(\mathrm{X}, \mu)\) 到 \(\mathcal{C}(\mathrm{Y})\) 中）是紧的。与典范映射 \(\mathcal{C}(\mathrm{X}) \to \mathcal{L}_{\mathbf{C}}^\infty(\mathrm{X}, \mu)\) 复合，便得到一个从 \(\mathcal{C}(\mathrm{X})\) 到 \(\mathcal{C}(\mathrm{Y})\) 中的紧线性映射。